\documentclass[11pt]{article}

% Packages
\usepackage[utf8]{inputenc}
\usepackage[margin=1in]{geometry}
\usepackage{amsmath, amssymb}
\usepackage{listings} % For code problems
\usepackage{xcolor}
\usepackage{graphicx}
% Code formatting
\lstset{
  basicstyle=\ttfamily\small,
  frame=single,
  breaklines=true,
  backgroundcolor=\color{gray!10}
}
\usepackage{hyperref}

\begin{document}

\begin{center}
    {\Large \textbf{AMATH 345: Practice Set \#1 (Fall 2026)} } \\
    \vspace{0.3cm}
   % \textit{Due at 6:00 PM, September $24^{th}$, 2025}
\end{center}

\hrule
\vspace{0.5em}

\emph{Note:} Marks assigned are meant to be representative of the relative weight that a similar question may be worth on a term test or exam.

\vspace{0.5em}
\hrule

\vspace{1cm}
\section*{Q1. (1 mark)}

Prove that \( A^TA\) shares the same null-space as \(A\). That is:
\[
\mathbf{N}(A) = \mathbf{N}(A^TA) 
\]

\section*{Q2. (4 marks)}


\begin{enumerate}
    \item Compute the eigenvalues and eigenvectors of the matrix \(A\). (2 mark)
\[
A=\begin{bmatrix}0 & 2 \\ 1 & 1\end{bmatrix},
\] 

    \item Verify numerically that:
    \begin{enumerate}
        \item The trace equals the sum of its eigenvalues. (0.5 mark)
        \item The determinant equals the product of its eigenvalues. (0.5 mark)
    \end{enumerate}

    \item Is A diagonalizable? If not, justify, if so, write A in the form \(A = X\Lambda X^{-1}\) (1 mark)
\end{enumerate}

\section*{Q3. (1 marks)}
 Suppose matrix \(X\) diagonalizes both matrices \(A\) and \(B\), and that they have the same eigenvectors. Then we  can write:
\[
A = X\Lambda_1 X^{-1} ,\quad B = X \Lambda_2 X^{-1}
\]
Prove that:
\[
AB = BA
\]

\section*{Q4. (3 marks)}

Let \(f:\mathbb{R}\rightarrow\mathbb{R}\) be convex and nondecreasing, and let
\(g:\mathbb{R}\rightarrow\mathbb{R}\) be convex.

Prove that the composition
\[
h(x) = f(g(x))
\]
is convex. Recall that:
\[
h:\mathbb{R}\rightarrow\mathbb{R}
\text{ is convex}
\]
if, for every \(\lambda\in[0,1]\) and any \(x,y\in\mathbb{R}\),
\[
h(\lambda x+(1-\lambda)y)
\leq
\lambda h(x)+(1-\lambda)h(y).
\]

\section*{Q5. (8 marks)}  %\texbf{Probably remove this!}}
In this question you will implement Newton's Method from scratch in Python and use it
to locate a local minimizer of the Rosenbrock function
\[
    f(x,y) = 100\,(y - x^2)^2 + (1 - x)^2.
\]

\begin{enumerate}
  \item \textbf{ Analytic gradient and Hessian.} (2 marks)\\
  Derive by hand the gradient and Hessian of $f(x,y)$. Show your work.

  \item \textbf{ Code the gradient and Hessian.} (1 marks)\\
  In Python, write functions that accept $(x,y)$ and return a NumPy vector and matrix:
\begin{verbatim}
def Grad(x, y):
    """Return numpy array [df/dx, df/dy] at (x,y)."""
    # your code

def Hessian(x, y):
    """Return 2x2 numpy array of second derivatives at (x,y)."""
    # your code
\end{verbatim}

  \item \textbf{ Implement multivariate Newton's method.} (3 marks)\\
  Write a function that uses your \texttt{Grad} and \texttt{Hessian} to iterate
  \[
    x_{k+1} = x_k + p_k, \quad \text{where } H(x_k)\,p_k = -\nabla f(x_k).
  \]
  Stop when $\|\nabla f(x_k)\|_2 < \texttt{tol}$ or when \texttt{max\_iter} is reached. 
  \emph{You may use \texttt{numpy.linalg} (e.g., \texttt{solve} or \texttt{inv}). Do not use optimization libraries.}
\newpage
\begin{verbatim}
def newton_method(f, grad, hess, x0, tol=1e-6, max_iter=50):
    """
    Apply Newton's method to minimize a two-variable function.

    Parameters
    ----------
    f : function
        Function f([x,y]) -> float. The Rosenbrock function in this problem.
    grad : function
        Gradient function grad(x,y) -> np.array([df/dx, df/dy]).
    hess : function
        Hessian function hess(x,y) -> 2x2 numpy array of second derivatives.
    x0 : array-like of shape (2,)
        Initial guess [x0, y0].
    tol : float, optional (default=1e-6)
        Stop if the 2-norm of the gradient is less than tol.
    max_iter : int, optional (default=50)
        Maximum number of Newton iterations.

    Returns
    -------
    x_star : np.array of shape (2,)
        Final point [x, y] after Newton iterations (approximate minimizer).
    history : list of dict
        A list containing information from each iteration. Each entry has keys:
            {
                "k"        : iteration number,
                "x"        : current point [x, y],
                "f"        : function value f(x,y),
                "grad_norm": ||grad f(x,y)||_2
            }

    Example
    -------
    >>> x_star, history = newton_method(f_rosenbrock, Grad, Hessian, [0,0])
    >>> history[0]
    {'k': 0, 'x': np.array([0.0, 0.0]), 'f': 1.0, 'grad_norm': 2.0}
    """
    # your code here
\end{verbatim}
\newpage
\item \textbf{Try (and plot!) multiple initial guesses and comment.} (2 marks) \\

Run your Newton's method implementation starting from at least three different 
initial guesses chosen from the domain
\[
x \in [-5,5], \quad y \in [-5,5].
\]

For each run, you are asked to:

\begin{enumerate}
    \item Plot the optimizer trajectory on a contour plot of the Rosenbrock function 
          (use filled contours and/or log-scale to clearly visualize the valley).
    \item Report the number of iterations until convergence 
          ($\|\nabla f(x_k)\|_2 < 10^{-6}$), or state \emph{``did not converge''} 
          if the method fails within the maximum iteration cap.
    \item Report the final point $(x^*, y^*)$ and the corresponding function value $f(x^*,y^*)$.
    \item Write a brief (1--2 sentence) comment on the observed behavior 
          (e.g., fast vs.\ slow convergence, sensitivity to the initialization).
\end{enumerate}

\textbf{Note:} You must select different initial guesses in the domain 
$[-5,5]\times[-5,5]$ and not simply start near the known minimizer $(1,1)$. 
The point of this exercise is to observe how Newton's method behaves globally.%, you will not be penalized for poor convergence.

\end{enumerate}

\textit{Constraints:} Do not use optimization libraries or automatic differentiation.
Use only NumPy linear algebra. Cap iterations at \texttt{max\_iter} and clearly
indicate if the method failed to meet the tolerance given your choice of initial guess.

\section*{Q6. (4 marks)}

Consider:
\[
A = 
\begin{bmatrix}
 1 & 2 \\ 2 & 4 \\ 1 & 2
\end{bmatrix}
\]

\begin{enumerate}
    \item find the eigenvalues and eigenvectors of \(A^TA\) (3 marks)
    \item Compute the "full" SVD of A. (1 mark)
\end{enumerate}

\section*{Q7. (2 marks)}
A symmetric matrix $S = S^\top$ has orthonormal eigenvectors $v_1, \dots, v_n$. Then any vector $x$ can be written as a combination
\[
x = c_1 v_1 + \cdots + c_n v_n. 
\]
Explain the following two formulas:
\[
x^\top x = c_1^2 + \cdots + c_n^2,
\quad \text{(1 mark)} \qquad 
x^\top Sx = \lambda_1 c_1^2 + \cdots + \lambda_n c_n^2. \quad \text{(1 mark)}
\]





\section*{Q8. (3 marks)}

\begin{enumerate}
    \item (1 mark)  Why can $r$ orthonormal vectors $u_{1}, \dots, u_{r}$ be   characterized by
\[
(m-1) + (m-2) + \cdots + (m-r)
\]
parameters?
    \item (2 marks) Show that an $m \times n$ matrix of rank $r$ has $r(m+n-r)$ free parameters in its singular value decomposition (SVD):
\[
A = U \Sigma V^{T} 
= \underbrace{(m \times r)}_{U} 
\underbrace{(r \times r)}_{\Sigma} 
\underbrace{(r \times n)}_{V^{T}}.
\]

\end{enumerate}

\section*{Q9. (4 marks)}

\[
\begin{aligned}
\text{minimize:} \quad & f(x_1,x_2) = x_1^2 + 2x_1x_2 + 2x_2^2 \\[6pt]
\text{subject to:} \quad & x_1 - x_2 = 2
\end{aligned}
\]


\section*{Q10. (6 marks)}

Let $A$ be an $m \times n$ real matrix and set $S = A^T A$.

\begin{enumerate}
    \item (2 marks) Show that maximizing $\|Ax\|^2$ over all nonzero vectors with \(\|x\| = 1 \) is equivalent to maximizing the Rayleigh quotient
    \[
        R(x) = \frac{x^T S x}{x^T x}
    \]
 over all nonzero vectors.
    \item (3 marks) Use the method of Lagrange multipliers to maximize $x^T S x$ subject to the constraint $\|x\|^2 = 1$. 
    Show that the maximizers must satisfy
    \[
        Sx = \lambda x
    \]
    for some scalar $\lambda$.

    \item (1 mark) Interpret the value of $\lambda$. What does it represent in terms of $\|Ax\|$?
\end{enumerate}

\newpage
\section*{Q11. (9 marks) (Programming)}

We work with the famous Iris dataset, which contains $150$ samples of flowers from three species
(\emph{setosa}, \emph{versicolor}, \emph{virginica}).  
You can either download the dataset or import it directly in Python from \href{https://archive.ics.uci.edu/dataset/53/iris}{here}.
Each sample has four measurements: sepal length, sepal width, petal length, and petal width.  Please answer all parts in your Jupyter Notebook, using markdown when necessary for comments or discussion.

\begin{figure}[htbp]
    \centering
    \includegraphics[width=0.7\linewidth]{Assessment&PracticePbs/PracticeSets/iris.png}
\end{figure}

Stack the data into a matrix
\[
X \in \mathbb{R}^{150 \times 4}.
\]

\begin{enumerate}
    \item (\textbf{Centering})  
    (1 mark) Compute the mean vector
    \[
    \mu = \frac{1}{150}\sum_{i=1}^{150} x_i,
    \]
    form the centered matrix $X_c = X - \mathbf{1}\mu^\top$, and display $\mu$.  
    Briefly interpret the meaning of $\mu$.

    \item (\textbf{SVD})  
    (2 marks) Compute the singular value decomposition
    \[
    X_c = U \Sigma V^\top.
    \]
    Interpret the right singular vectors (columns of $V$) in this context.  
    Which flower measurements (sepal length, sepal width, petal length, petal width) dominate the first two singular vectors?

    \item (\textbf{Projection}) (4 marks)  
    Project the data onto the leading singular vectors and plot the $150$ points in $\mathbb{R}^k$ (color coded by species):
    \begin{itemize}
        \item \(k=1\): projection onto the first singular vector.
        \item \(k=2\): projection onto the first two singular vectors.
        \item \(k=3\): projection onto the first three singular vectors.
    \end{itemize}
    Comment on how the species appear in these lower-dimensional visualizations.

    \item (\textbf{Discussion}) (2 marks)
    \begin{itemize}
        \item Why do the first two singular vectors already give a good visualization of the dataset?  
        \item What do the singular values represent?
    \end{itemize}
\end{enumerate}


\end{document}

\end{document}

